\section*{الفصل \ref{ch:b1:lab-techniques-1} --- التقنيات المخبرية (1): السلامة والقياس والفصل}

\begin{solution}{exo:b1:lab-techniques-1:1}
\emph{خطر}: يستدعيها H225 (سائل شديد الاشتعال، الفئة 2)،
وتحمل البطاقة الكلمة الأشد. وH225 \omterm{def:b1:lab-techniques-1:hazard}{خطر} فيزيائي؛ أما H319
(تهيج شديد للعين) وH336 (نعاس أو دوار) فخطران
صحيان. الاحتياطات: لا لهب ولا شرر بالقرب، وقارورة مغلقة؛ ونظارات؛
والعمل في مكان مهوّى أو في خزانة غازات.
% ledger: ghs:acetone, hstat:H225, hstat:H319, hstat:H336
\end{solution}

\begin{solution}{exo:b1:lab-techniques-1:2}
(أ) \omterm{def:b1:lab-techniques-1:hazard}{خطر}، وهو خاصية للحمض. (ب) \omterm{def:b1:lab-techniques-1:hazard}{خطورة}، صغيرة لأن
التعرّض صغير. (ج) \omterm{def:b1:lab-techniques-1:hazard}{خطر}. (د) \omterm{def:b1:lab-techniques-1:hazard}{خطورة}: \omterm{def:b1:lab-techniques-1:hazard}{الخطر} نفسه، والتعرّض أكبر
(البخار، والرذاذ) حين يكون ساخنًا ومفتوحًا.
\end{solution}

\begin{solution}{exo:b1:lab-techniques-1:3}
$R_f = 2.1/6.0 = 0.35$ و$4.2/6.0 = 0.70$. للبقعة الثانية $R_f$
المرجع نفسه: فقد تحتويه العيّنة (وهذا متسق)، لكن
تساوي $R_f$ \omterm{def:b1:extent-q-and-k:system}{بطور} متحرك واحد لا يثبت التطابق؛ \omterm{def:b1:extent-q-and-k:system}{وطور} متحرك ثانٍ، أو
بقعة مشتركة من العيّنة والمرجع تبقى بقعة واحدة، يعزّز
الاستنتاج. أما البقعة عند 0.35 فهي يقينًا نوع كيميائي آخر.
\end{solution}

\begin{solution}{exo:b1:lab-techniques-1:4}
نصف العرض \qty{0.05}{mL}، توزيع مستطيل: $u = 0.05/\sqrt3 = \qty{0.029}{mL}$.
ومن أجل فرق بين قراءتين مستقلتين،
$u = \sqrt{2} \times 0.029 = \qty{0.041}{mL}$.
\end{solution}

\begin{solution}{exo:b1:lab-techniques-1:5}
المتوسط \qty{0.25100}{g}؛ والانحرافات $+2$، $-2$، $+5$، 0، $-5$ (بالوحدة
\qty{e-4}{g})، ومجموع المربعات \qty{58e-8}{g^2}، $s = \sqrt{58 \times
10^{-8}/4} = \qty{3.8e-4}{g}$؛ $u = s/\sqrt5 = \qty{1.7e-4}{g}$؛
$U = \qty{3.4e-4}{g}$: \babelsublr{$m = (0.25100 \pm 0.00034)$~g}، $k = 2$.
\end{solution}

\begin{solution}{exo:b1:lab-techniques-1:6}
بجزء واحد: $f = 1/(1 + 3 \times 30/50) = 0.357$، أي \qty{64.3}{\percent}
مستخلصًا. وبجزأين من \qty{15}{mL}: $f = 1.9^{-2} = 0.277$،
أي \qty{72.3}{\percent}. وبثلاثة من \qty{10}{mL}: $f = 1.6^{-3} = 0.244$،
أي \qty{75.6}{\percent}.
\end{solution}

\begin{solution}{exo:b1:lab-techniques-1:7}
$z = 0.0012/\sqrt{0.0006^2 + 0.0002^2} = 0.0012/0.00063 = 1.9 \le 2$:
متوافقتان، بالكاد. ومع \qty{0.0004}{mol/L}: $z = 0.0012/0.00045 = 2.7$،
غير متوافقتين: \omterm{def:b1:titration-methods:titration}{فالمعايرة} الأدق تكشف فرقًا حقيقيًا
(خطأ في التحضير، أو في \omterm{def:b1:titration-methods:titration}{المعايرة}).
\end{solution}

\begin{solution}{exo:b1:lab-techniques-1:8}
Q: قليل الذوبان باردًا، شديد الذوبان ساخنًا. أما P فلا يذيب الصلب ساخنًا أفضل
بكثير منه باردًا (فلا يتبلور شيء عند التبريد)؛ وR يُبقي الكثير
ذائبًا باردًا (خسائر كبيرة). ومع Q، يحتاج \qty{5.0}{g} إلى ما لا يقل عن
$5.0/12 \times 100 = \qty{42}{mL}$ من المذيب المغلي، تُبقي
$0.3 \times 0.417 = \qty{0.13}{g}$ ذائبًا باردًا: فيمكن استرجاع
\qty{4.9}{g} على الأكثر، أي \qty{97.5}{\percent} (قبل أي فقد على
المرشح).
\end{solution}

\begin{solution}{exo:b1:lab-techniques-1:9}
الإيثانول ($n_D = 1.3611$ عند \qty{20}{\celsius})؛ فالماء (1.333)
والهكسان الحلقي (1.4266) مستبعدان. ويتغير المعامل مع درجة الحرارة،
فلا معنى للقيمة دونها. أما القيمة 1.350، الواقعة بين الماء
والإيثانول، فتوحي بخليط، كإيثانول يحتوي ماء.
% ledger: nD:water, nD:ethanol, nD:cyclohexane
\end{solution}

\begin{solution}{exo:b1:lab-techniques-1:10}
مع $t = KV/V_a$ و$n$ جزءًا، $\ln f_n = -n\ln(1 + t/n)$؛ وحين
$n \to \infty$، يكون $n \ln(1 + t/n) \to t$ (لأن $\ln(1 + \varepsilon) \sim
\varepsilon$)، و$f_n$ تتناقص نحو $\mathrm{e}^{-t}$ (\cref{prop:b1:lab-techniques-1:multiple-extractions}).
وهنا $t = 3 \times 30/50 = 1.8$، $\mathrm{e}^{-1.8} = 0.165$: أي
\qty{83.5}{\percent} على الأكثر. وثلاثة أجزاء تعطي من قبل \qty{75.6}{\percent}،
أي \qty{91}{\percent} من النهاية؛ وبلوغ \qty{99}{\percent} منها يتطلب 32
جزءًا حجم كل منها أقل من \qty{1}{mL}: فبعد جزأين أو ثلاثة، لا يستحق الكسب
العناء.
\end{solution}

\begin{solution}{exo:b1:lab-techniques-1:11}
مع $t = x - x_0$: $\int_{-a}^{a} (a - |t|)/a^2 \,\mathrm{d}t =
2 \int_0^a (a - t)/a^2 \,\mathrm{d}t = 2 (a^2/2)/a^2 = 1$. والمتوسط
$x_0$ بالتناظر؛ والتباين $2 \int_0^a t^2 (a - t)/a^2
\,\mathrm{d}t = \frac{2}{a^2}\left(\frac{a^4}{3} - \frac{a^4}{4}\right) =
\frac{a^2}{6}$، ومن ثم $u = a/\sqrt6 = 0.41\,a$، مقابل $a/\sqrt3 = 0.58\,a$
في الحالة المستطيلة: فمعرفة أن الوسط أرجح تخفض
\omterm{def:b1:lab-techniques-1:uncertainty}{عدم اليقين}.
\end{solution}

\begin{solution}{exo:b1:lab-techniques-1:12}
$c = m/(MV) = 0.25100/(204.2 \times 0.10000) = \qty{0.012292}{mol/L}$.
أوجه \omterm{def:b1:lab-techniques-1:uncertainty}{عدم اليقين} النسبية: الكتلة $1.7 \times 10^{-4}/0.2510 = 6.8 \times
10^{-4}$، والحجم $0.06/100.00 = 6.0 \times 10^{-4}$؛ والمركّب
$\sqrt{(6.8^2 + 6.0^2)} \times 10^{-4} = 9.1 \times 10^{-4}$، أي
\qty{0.091}{\percent}؛ $u(c) = \qty{1.1e-5}{mol/L}$، $U = \qty{2.2e-5}{mol/L}$:
\babelsublr{$c = (0.012292 \pm 0.000022)$~mol/L}، $k = 2$. ويطغى الوزن،
قليلًا؛ فالمصدران من الحجم نفسه.
\end{solution}

\begin{solution}{pb:b1:lab-techniques-1:1}
\textbf{1.} GHS02 قابل للاشتعال، GHS05 أكّال، GHS06 سمية حادة، GHS07
ضار أو مهيّج.
\textbf{2.} حمض الساليسيليك: \emph{خطر}، بسبب H318 (ضرر شديد
للعين، الفئة 1). الأسبرين: \emph{تحذير} (H302 وحده).
\textbf{3.} \omterm{def:b1:lab-techniques-1:hazard}{الخطر} ثابت، أما التعرّض فلا: حجم صغير،
يُقاس ويُستعمل في خزانة غازات، مع القفازات والنظارات والمئزر،
والقارورة تُغلق فورًا، ولا لهب بالقرب.
\textbf{4.} H314 (حروق جلدية شديدة وضرر للعين): النظارات والقفازات.
وH332 (ضار إذا استُنشق) وH226 (سائل وبخار قابلان للاشتعال): خزانة
الغازات وغياب اللهب.
\textbf{5.} اشطف بحذر بالماء لعدة دقائق، مع نزع
العدسات اللاصقة إن أمكن، وواصل الشطف (\babelsublr{P305+P351+P338})؛ ثم اطلب
المشورة الطبية.
\textbf{6.} إلى النفايات المائية، بعد المعادلة، لا إلى
المغسلة.
\textbf{7.} \ce{C7H6O3 + C4H6O3 -> C9H8O4 + C2H4O2}: C $7 + 4 = 9 + 2$،
H $6 + 6 = 8 + 4$، O $3 + 3 = 4 + 2$.
\textbf{8.} $2.00/138.0 = \qty{1.449e-2}{mol}$؛ على الأكثر
$\num{1.449e-2} \times 180.0 = \qty{2.61}{g}$ من الأسبرين.
\textbf{9.} يحتوي الصلب الخام على حمض ساليسيليك لم يتفاعل، وحمض
الإيثانويك والماء: فكتلته ليست كتلة الأسبرين، ونقاوته
مجهولة.
\textbf{10.} ما يبقى ذائبًا بعد التبريد يُفقد: \omterm{def:b1:precipitation:solubility}{فالذوبانية} الصغيرة
باردًا والكبيرة ساخنًا تعني أن الناتج يذوب ساخنًا ويعود
كله تقريبًا باردًا.
\textbf{11.} كل مليلتر إضافي يُبقي حصته من الناتج ذائبة
باردًا.
\textbf{12.} النمو البطيء يعطي بلورات أكبر وأكثر انتظامًا تحتجز
قدرًا أقل من السائل الأم والشوائب؛ والثلج يخفض \omterm{def:b1:precipitation:solubility}{الذوبانية}، ومن ثم
الفقد.
\textbf{13.} $\qty{0.045}{L} \times \qty{3.3}{g/L} = \qty{0.15}{g}$.
\textbf{14.} $1.95/2.61 = \qty{75}{\percent}$ (\qty{74.7}{\percent}).
\textbf{15.} $807/6 = \qty{134.5}{\celsius}$.
\textbf{16.} الانحرافات $-0.5$، $0.5$، $-1.5$، $0.5$، $-0.5$، $1.5$؛ ومجموع
المربعات 5.5؛ $s = \sqrt{5.5/5} = \qty{1.05}{\celsius}$.
\textbf{17.} $u_A = 1.05/\sqrt6 = \qty{0.43}{\celsius}$.
\textbf{18.} نصف العرض \qty{0.5}{\celsius}: $0.5/\sqrt3 = \qty{0.29}{\celsius}$.
\textbf{19.} $1/\sqrt3 = \qty{0.58}{\celsius}$؛
$u = \sqrt{0.428^2 + 0.289^2 + 0.577^2} = \sqrt{0.600} = \qty{0.77}{\celsius}$.
\textbf{20.} $U = 2 \times 0.775 = \qty{1.5}{\celsius}$:
\babelsublr{$T = (134.5 \pm 1.5)$~\unit{\celsius}}، $k = 2$.
\textbf{21.} ينصهر عند درجة منخفضة وعلى مجال قدره \qty{8}{\celsius}: فهو
غير نقي (حمض الساليسيليك، وحمض الإيثانويك، والماء).
\textbf{22.} لكونها معطاة إلى الوحدة، تقع القيمة ضمن $\pm\qty{0.5}{\celsius}$:
$u_{\mathrm{ref}} = 0.5/\sqrt3 = \qty{0.29}{\celsius}$.
\textbf{23.} $2.3/6.0 = 0.38$ و$3.3/6.0 = 0.55$. فالناتج الخام
كان يحتوي حمض الساليسيليك (القيمة $R_f$ نفسها للمرجع) والأسبرين؛ وبعد
\omterm{def:b1:lab-techniques-1:recrystallisation}{إعادة التبلور} لا تبقى إلا بقعة الأسبرين: فقد أُزيل حمض الساليسيليك،
على الأقل إلى ما دون ما تكشفه الصفيحة.
\textbf{24.} يتفكك الأسبرين أثناء انصهاره؛ وإذا سُخّن ببطء، كان لديه وقت كي
يتفكك، فينصهر الخليط المتكوّن عند درجة أدنى. والقيمة المجدولة تخص
التسخين السريع، كما على المنصة.
\textbf{25.} $z = |134.5 - 135|/\sqrt{0.600 + 0.083} = 0.5/0.827 =$
\textbf{0.60}: أدنى بكثير من 2، فدرجتا الانصهار المقيسة والمجدولة
متوافقتان.
\end{solution}
